\documentclass[12pt,a4paper]{article}

\usepackage[utf8]{inputenc}
\usepackage[T1]{fontenc}
\usepackage[brazil]{babel}
\usepackage{geometry}
\usepackage{amsmath}
\usepackage{hyperref}
\usepackage{enumitem}
\usepackage{booktabs}
\usepackage{array}
\usepackage{parskip}
\usepackage{titlesec}
\usepackage{fancyhdr}
\usepackage{amssymb}

\geometry{
  top=2.5cm,
  bottom=2.5cm,
  left=3cm,
  right=2.5cm
}

% Header/footer
\pagestyle{fancy}
\fancyhf{}
\rhead{CTJ 2026-1 EMB5476}
\cfoot{\thepage}
\renewcommand{\headrulewidth}{0.4pt}

% Section formatting
\titleformat{\section}{\large\bfseries}{\thesection}{1em}{}
\titleformat{\subsection}{\normalsize\bfseries}{\thesubsection}{1em}{}

% Hyperref setup
\hypersetup{
  colorlinks=true,
  urlcolor=blue,
  linkcolor=black
}

% Checkbox command
\newcommand{\checkbox}{$\square$}

\begin{document}

% Title block
\begin{center}
  {\large\bfseries CTJ 2026-1 \quad EMB5476}\\[0.5em]
  {\large\bfseries Trabalho 1 -- Tutoriais GT1, GT2 e GT3}\\[0.3em]
  {\normalsize Entrega limite até as 22 horas do dia 15/06/26}\\[0.2em]
  {\normalsize Grupos de até 4 alunos}
\end{center}

\vspace{1em}
\hrule
\vspace{1em}

% ─────────────────────────────────────────────
\section{Objetivos}

Este trabalho tem como objetivos:

\begin{enumerate}
  \item Comparar diferentes abordagens de modelagem do escoamento (laminar, RANS e modelos de transição).
  \item Avaliar o impacto do tratamento de parede (\textit{wall-function} vs.\ \textit{wall-resolved}) na qualidade dos resultados.
  \item Executar e interpretar rotinas de pós-processamento nativas dos casos (\texttt{gnuplot}, \texttt{foamPostProcess}, amostragens).
  \item Desenvolver pelo menos um pós-processamento próprio (script em Python/Matlab/gnuplot/bash).
  \item Sustentar tecnicamente as escolhas de malha, condições de contorno (BCs) e modelo com base em teoria e em métricas (perfil de velocidade, $C_f$, $y^+$, etc.).
\end{enumerate}

% ─────────────────────────────────────────────
\section{Casos e Escopo}

\subsection*{GT1 -- Escoamento em duto (laminar e turbulento)}

Diretórios recomendados:
\begin{itemize}
  \item \texttt{GT1/laminar}
  \item \texttt{GT1/turbulent\_wedge} (ou \texttt{turbulent\_planar}, justificar escolha)
\end{itemize}

\subsection*{GT2 -- Placa plana (laminar, turbulento e variações de modelos)}

Diretórios recomendados (escolher um conjunto representativo):
\begin{itemize}
  \item \texttt{GT2/RANS/KEpsilon\_highRE\_V1}
  \item \texttt{GT2/RANS/KOmegaSST\_highRE}
  \item \texttt{GT2/RANS/SA\_LRN} (ou outro low-Re)
  \item \texttt{GT2/SRS/DES\_KO\_LRN} (comparação SRS)
  \item \texttt{GT2/SRS/LES-WALE\_LRN} (comparação SRS)
\end{itemize}

\textbf{Observação:} é obrigatório comparar pelo menos um caso high-Re (com \textit{wall functions}) e um low-Re/\textit{wall-resolved}. Também é obrigatório realizar uma comparação específica entre as abordagens 2D \texttt{DES\_KO\_LRN} e \texttt{LES-WALE\_LRN}, discutindo limitações físicas dessa aproximação em 2D.

\subsection*{GT3 -- Placa plana com transição (ERCOFTAC T3A)}

Diretórios recomendados:
\begin{itemize}
  \item \texttt{GT3/RANS-kklo}
  \item \texttt{GT3/RANS-kOmegaSSTLM\_UzG} (ou \texttt{RANS-kOmegaSSTLM})
  \item \texttt{GT3/setup\_TM/...} para \textit{benchmark} com modelos totalmente turbulentos
\end{itemize}

% ─────────────────────────────────────────────
\section{Procedimento Mínimo Obrigatório}

\subsection{Etapa A -- Inspeção inicial do caso}

Para cada caso executado:

\begin{enumerate}[label=\alph*)]
  \item Ler \texttt{README*} e scripts \texttt{run\_all.sh}, \texttt{run\_mesh.sh}, \texttt{run\_solver.sh}.
  \item Identificar no relatório:
  \begin{itemize}
    \item solver utilizado (\texttt{foamRun}/\texttt{incompressibleFluid});
    \item sequência de execução real;
    \item se há execução serial ou paralela;
    \item campos turbulentos adicionais requeridos ($k$, $\omega$, $\epsilon$, $\nu_t$, etc.).
  \end{itemize}
\end{enumerate}

\subsection{Etapa B -- Execução}

Executar por caso, no mínimo:

\begin{enumerate}[label=\alph*)]
  \item \textbf{Malha:} \texttt{sh run\_mesh.sh} e validação via \texttt{log.checkMesh}.
  \item \textbf{Solução:} \texttt{sh run\_solver.sh} (ou \texttt{sh run\_all.sh} quando incluir pós-processamento).
  \item \textbf{Pós-processamento:} executar scripts do caso (\texttt{run\_gnuplot.sh}, etc.) e comandos \texttt{foamPostProcess} relevantes.
\end{enumerate}

\subsection{Etapa C -- Variação de configuração (estudo)}

Obrigatório realizar alterações controladas (uma por vez), por exemplo:
\begin{itemize}
  \item troca de modelo em \texttt{constant/momentumTransport};
  \item troca de tratamento de parede em \texttt{0/nut}, \texttt{0/k}, \texttt{0/omega}, \texttt{0/epsilon};
  \item ajuste de intensidade de turbulência de entrada e/ou razão $\mu_t/\mu$;
  \item refinamento local de malha perto da parede (quando aplicável).
\end{itemize}

Cada alteração deve ser acompanhada por uma explicação.

% ─────────────────────────────────────────────
\section{Análises Técnicas Esperadas}

\subsection{GT1 -- Duto}

\begin{itemize}
  \item Comparar perfil de velocidade laminar numérico com perfil analítico de Poiseuille.
  \item Comparar perfil turbulento com lei de potência/\textit{log-law} (comentar limitações).
  \item Calcular/avaliar $y^+$ e tensão cisalhante na parede (\texttt{wallShearStress}).
  \item Discutir efeito da escolha de \textit{wedge} vs.\ \textit{planar} (custo e fidelidade).
\end{itemize}

\subsection{GT2 -- Placa plana}

\begin{itemize}
  \item Construir gráfico $u^+$ vs.\ $y^+$ para ao menos 2 modelos.
  \item Calcular coeficiente de atrito local:
  \begin{equation}
    C_f = \frac{\tau_w}{\tfrac{1}{2}\,\rho\,U_\infty^2}
  \end{equation}
  e comparar tendência com correlações/clássico da placa plana.
  \item Comparar resposta high-Re (\textit{wall functions}) vs.\ low-Re (\textit{wall-resolved}).
  \item Comparar explicitamente os casos \texttt{GT2/SRS/DES\_KO\_LRN} e \texttt{GT2/SRS/LES-WALE\_LRN} com base em $u^+$-$y^+$, $C_f(x)$, $y^+$, custo computacional e estabilidade numérica.
  \item Comentar criticamente as limitações de usar DES e LES em 2D (conforme notas dos próprios casos) e o impacto disso na interpretação dos resultados.
  \item Relacionar qualidade de resultados com faixa de $y^+$ observada.
\end{itemize}

\subsection{GT3 -- Transição (T3A)}

\begin{itemize}
  \item Comparar pelo menos um modelo de transição (ex.: $k$-$k_L$-$\omega$ ou SST-LM) contra um modelo totalmente turbulento.
  \item Identificar ponto aproximado de transição (via $C_f(x)$, crescimento de $\nu_t$, ou mudança do perfil).
  \item Discutir sensibilidade a condições de entrada ($\mathrm{TU}$, escalas turbulentas).
  \item Explicar porque modelos de transição tendem a exigir abordagem \textit{wall-resolved}.
\end{itemize}

% ─────────────────────────────────────────────
\section{Pós-processamento Próprio (obrigatório)}

Cada grupo deve criar ao menos um script próprio (\texttt{.py}, \texttt{.ipynb}, \texttt{.m}, ou \texttt{.sh}) para automatizar uma análise. Exemplos:
\begin{itemize}
  \item cálculo automático de $u^+$ e $y^+$ a partir de amostras;
  \item comparação de $C_f(x)$ entre 3 modelos;
  \item extração de métricas de convergência de \texttt{log.solver}.
\end{itemize}

O script deve ser entregue com:
\begin{itemize}
  \item instruções de execução;
  \item entradas esperadas;
  \item saídas geradas (figuras/tabela).
\end{itemize}

% ─────────────────────────────────────────────
\section{Estrutura Obrigatória do Relatório}

\begin{enumerate}
  \item \textbf{Introdução} (contexto: turbulência, fechamento, parede, transição).
  \item \textbf{Metodologia} (casos, malhas, BCs, modelos, comandos usados).
  \item \textbf{Resultados} (figuras comparativas, tabelas, métricas).
  \item \textbf{Discussão} (causa física das diferenças; não apenas ``funcionou/não funcionou'').
  \item \textbf{Conclusões} (o que aprenderam sobre modelo vs.\ malha vs.\ parede).
  \item \textbf{Anexo} com scripts e comandos.
\end{enumerate}

% ─────────────────────────────────────────────
\section{Checklist de Qualidade Mínima}

\begin{itemize}[label=\checkbox]
  \item Mostrar evidências de \texttt{checkMesh} e convergência/resíduos.
  \item Reportar faixa de $y^+$ nos casos turbulentos.
  \item Comparar pelo menos 3 configurações de turbulência no total (somando GT1+GT2+GT3).
  \item Incluir no GT2 a comparação \texttt{DES\_KO\_LRN} vs.\ \texttt{LES-WALE\_LRN} com discussão das limitações de SRS em 2D.
  \item Incluir ao menos 1 comparação com referência analítica/correlação/benchmark.
  \item Incluir script próprio de pós-processamento.
  \item Discutir custo computacional (tempo, iterações, paralelo vs.\ serial quando aplicável).
\end{itemize}

% ─────────────────────────────────────────────
\section{Critérios de Avaliação}

\begin{center}
\begin{tabular}{>{\raggedright\arraybackslash}p{10cm} r}
  \toprule
  \textbf{Item} & \textbf{Peso} \\
  \midrule
  Fundamentação física (turbulência, parede, transição)          & 25\% \\
  Qualidade técnica das simulações (setup, consistência, validação) & 25\% \\
  Análise crítica comparativa (não apenas descritiva)             & 25\% \\
  Pós-processamento próprio e reprodutibilidade                   & 15\% \\
  Organização, clareza e qualidade das figuras/tabelas            & 10\% \\
  \bottomrule
\end{tabular}
\end{center}

% ─────────────────────────────────────────────
\section{Referências Recomendadas}

\begin{itemize}
  \item Material de apoio do curso: \textit{Turbulence modeling in OpenFOAM: Guided Tutorials} (seção avançada, 2025).
  \item NASA Turbulence Modeling Resource: \url{https://turbmodels.larc.nasa.gov/}
  \item ERCOFTAC T3A (caso transicional): links fornecidos no material.
\end{itemize}

\vspace{1em}
\noindent\textbf{Observação final:} a ênfase deste trabalho é demonstrar entendimento, não apenas executar scripts. Resultados ``bons'' sem explicação física consistente terão nota limitada.

\newpage


\section{Introdução}


\section{Casos}

    \begin{enumerate}
    
        \item GT1

            \begin{itemize}
                
                \item Laminar -  listar (Fazer em tabela unificada?)

                     \begin{itemize}
                        \item solver utilizado (\texttt{foamRun}/\texttt{incompressibleFluid});
                        \item sequência de execução real;
                        \item se há execução serial ou paralela;
                        \item campos turbulentos adicionais requeridos ($k$, $\omega$, $\epsilon$, $\nu_t$, etc.).
                      \end{itemize}

                        Estudo de caso - alterar um por vez = 

                    \begin{itemize}
                      \item troca de modelo em \texttt{constant/momentumTransport};
                      \item troca de tratamento de parede em \texttt{0/nut}, \texttt{0/k}, \texttt{0/omega}, \texttt{0/epsilon};
                      \item ajuste de intensidade de turbulência de entrada e/ou razão $\mu_t/\mu$;
                      \item refinamento local de malha perto da parede (quando aplicável).
                    \end{itemize}
                    
                \item Turbulento - GT1/turbulent wedge | GT1/turbulent\_planar - Justificar escolha

                     \begin{itemize}
                        \item solver utilizado (\texttt{foamRun}/\texttt{incompressibleFluid});
                        \item sequência de execução real;
                        \item se há execução serial ou paralela;
                        \item campos turbulentos adicionais requeridos ($k$, $\omega$, $\epsilon$, $\nu_t$, etc.).
                      \end{itemize}

                        Estudo de caso - alterar um por vez = 

                    \begin{itemize}
                      \item troca de modelo em \texttt{constant/momentumTransport};
                      \item troca de tratamento de parede em \texttt{0/nut}, \texttt{0/k}, \texttt{0/omega}, \texttt{0/epsilon};
                      \item ajuste de intensidade de turbulência de entrada e/ou razão $\mu_t/\mu$;
                      \item refinamento local de malha perto da parede (quando aplicável).
                    \end{itemize}
                
            \end{itemize}

            \begin{enumerate}
                
                \item Comparar perfil de velocidade laminar numérico com perfil analítico de Poiseuille. 

                \item Comparar perfil turbulento com lei de potência/log-law (comentar limitações).

                \item Calcular/avaliar y+ e tensão cisalhante na parede (wallShearStress).

                \item Discutir efeito da escolha de wedge vs. planar (custo e fidelidade).
                
            \end{enumerate}

        \item GT2 

            \begin{itemize}
                
                \item Turbulento High Re

                \item Turbulento Low Re

                \item comparação entre DES KO LRN e LES-WALE LRN, discutindo limitações fisicas dessa aproximaçãoo em 2D.
                
            \end{itemize}

            \begin{enumerate}
                
                \item Construir grafico $u^+$ vs. $y^+$ para ao menos 2 modelos.

                \item Calcular coeficiente de atrito local e comparar tendência com correlações/clássico da placa plana.
                
                    \begin{equation}
                        C_f = \frac{\tau_w}{\tfrac{1}{2} \rho U_\infty^2}
                        \label{eq:placeholder_label}
                    \end{equation}

                \item Comparar resposta high-Re (wall functions) vs. low-Re (wall-resolved)

                \item  Comparar explicitamente os casos GT2/SRS/DES KO LRN e GT2/SRS/LES-WALE LRN com base em mu+-y+, Cf (x), y+, custo computacional e estabilidade numérica

                \item Comentar criticamente as limitações de usar DES e LES em 2D ( conforme notas dos próprios casos) e o impacto disso na interpretação dos resultados.

                \item Relacionar qualidade de resultados com faixa de y+ observada.

            \end{enumerate}

        \item GT3

            \begin{itemize}
                
                \item GT3/RANS-kklo

                \item GT3/RANS-kOmegaSSTLM

                \item GT3/setup TM/... para benchmark com modelos totalmente turbulentos
                
            \end{itemize}

            \begin{enumerate}

                \item Comparar pelo menos um modelo de transição (ex.: k-kl-$\omega$ ou SST-LM) contra um modelo totalmente turbulento.

                \item Identificar ponto aproximado de transição (via Cf (x), crescimento de $\nu_t$, ou mudan\c{c}a do perfil).

                \item Discutir sensibilidade a condições de entrada (TU, escalas turbulentas)

                \item Explicar porque modelos de transição tendem a exigir abordagem wall-resolved.
                
            \end{enumerate}
    
    \end{enumerate}

\end{document}
